\nonstopmode{}
\documentclass[a4paper]{book}
\usepackage[times,hyper]{Rd}
\usepackage{makeidx}
\makeatletter\@ifl@t@r\fmtversion{2018/04/01}{}{\usepackage[utf8]{inputenc}}\makeatother
% \usepackage{graphicx} % @USE GRAPHICX@
\makeindex{}
\begin{document}
\chapter*{}
\begin{center}
{\textbf{\huge Package `lmmprobe'}}
\par\bigskip{\large \today}
\end{center}
\ifthenelse{\boolean{Rd@use@hyper}}{\hypersetup{pdftitle = {lmmprobe: Sparse High-Dimensional Linear Mixed Modeling with a Partitioned Empirical Bayes ECM Algorithm}}}{}
\ifthenelse{\boolean{Rd@use@hyper}}{\hypersetup{pdfauthor = {Anja Zgodic; Alex McLain; Peter Olejua}}}{}
\begin{description}
\raggedright{}
\item[Type]\AsIs{Package}
\item[Title]\AsIs{Sparse High-Dimensional Linear Mixed Modeling with a Partitioned
Empirical Bayes ECM Algorithm}
\item[Version]\AsIs{1.0.1}
\item[Date]\AsIs{2026-02-18}
\item[Description]\AsIs{Implements a partitioned Empirical Bayes ECM algorithm for sparse high-dimensional linear mixed modeling. The package provides efficient estimation and inference for mixed models with high-dimensional fixed effects.}
\item[License]\AsIs{GPL (>= 2)}
\item[Encoding]\AsIs{UTF-8}
\item[RoxygenNote]\AsIs{7.3.3}
\item[Depends]\AsIs{R (>= 3.5.0)}
\item[Imports]\AsIs{Rcpp (>= 1.0.8.3), tidyr (>= 1.2.0), lme4 (>= 1.1-29), future
(>= 1.28.0), future.apply (>= 1.10.0)}
\item[LinkingTo]\AsIs{Rcpp, RcppArmadillo}
\item[Suggests]\AsIs{testthat (>= 3.0.0), knitr, rmarkdown, MASS}
\item[VignetteBuilder]\AsIs{knitr}
\item[NeedsCompilation]\AsIs{yes}
\item[Author]\AsIs{Anja Zgodic [aut, cre],
Alex McLain [aut],
Peter Olejua [aut]}
\item[Maintainer]\AsIs{Anja Zgodic }\email{azgodic@email.sc.edu}\AsIs{}
\end{description}
\Rdcontents{Contents}
\HeaderA{lmmprobe-package}{lmmprobe: Sparse High-Dimensional Linear Mixed Modeling with a Partitioned Empirical Bayes ECM Algorithm}{lmmprobe.Rdash.package}
\keyword{internal}{lmmprobe-package}
%
\begin{Description}
Implements a partitioned Empirical Bayes ECM algorithm for sparse high-dimensional linear mixed modeling. The package provides efficient estimation and inference for mixed models with high-dimensional fixed effects.
\end{Description}
%
\begin{Author}
\strong{Maintainer}: Anja Zgodic \email{azgodic@email.sc.edu}

Authors:
\begin{itemize}

\item{} Alex McLain
\item{} Peter Olejua

\end{itemize}


\end{Author}
\HeaderA{lmmprobe}{Sparse high-dimensional linear mixed modeling with PaRtitiOned empirical Bayes ECM (LMM-PROBE) algorithm.}{lmmprobe}
%
\begin{Description}
Sparse high-dimensional linear mixed modeling with PaRtitiOned empirical Bayes ECM (LMM-PROBE) algorithm. Currently, the package offers functionality for two scenarios. Scenario 1: only a random intercept, each unit has the same number of observations; Scenario 2: a random intercept and a random slope, each unit has the same number of observations. We are actively expanding the package for more flexibility and scenarios.
\end{Description}
%
\begin{Arguments}
\begin{ldescription}
\item[\code{Y}] A training-data matrix containing the outcome \code{Y}.

\item[\code{Z}] A training-data matrix containing the sparse fixed-effect predictors on which to apply the lmmprobe algorithm. The first columns should be the "id" column.

\item[\code{V}] A training-data matrix containing non-sparse predictors for the random effects. This matrix is currently only programmed for two scenarios. Scenario 1: only a random intercept, where V is a matrix with one column containing ID's and each unit has the same number of observations. Scenario 2: a random intercept and a random slope, where V is a matrix with two columns. The first column is ID and the second column is a continuous variable (e.g. time) for which a random slope is to be estimated. Each unit has the same number of observations.

\item[\code{ID\_data}] A factor vector of IDs for subjects in the training set.

\item[\code{Y\_test}] A testing-data matrix containing the outcome \code{Y}. Default is NULL.

\item[\code{Z\_test}] A training-data matrix containing the sparse fixed-effect predictors on which to apply the lmmprobe algorithm. The first columns should be the "id" column. Default is NULL.

\item[\code{V\_test}] A training-data matrix containing non-sparse predictors for the random effects. This matrix is currently only programmed for two scenarios. Scenario 1: only a random intercept, where V is a matrix with one column containing ID's and each unit has the same number of observations. Scenario 2: a random intercept and a random slope, where V is a matrix with two columns. The first column is ID and the second column is a continuous variable (e.g. time) for which a random slope is to be estimated. Each unit has the same number of observations. Default is NULL.

\item[\code{ID\_test}] A factor vector of IDs for subjects in the testing set.

\item[\code{alpha}] Type I error; significance level.

\item[\code{ep}] Value against which to compare convergence criterion, we recommend 0.05.

\item[\code{B}] The number of groups to categorize estimated coefficients in to calculate correlation \eqn{\rho}{}. We recommend five.

\item[\code{adj}] A factor multiplying Silverman’s 'rule of thumb' in determining the bandwidth for density estimation, same as the 'adjust' argument of R's density function. Default is three.

\item[\code{maxit}] Maximum number of iterations the algorithm will run for. Default is 10000.

\item[\code{cpus}] The number of CPUS user would like to use for parallel computations. Default is four.

\item[\code{LR}] A learning rate parameter \code{r}. Using zero corresponds to the implementation described in Zgodic et al.

\item[\code{C}] A learning rate parameter \code{c}. Using one corresponds to the implementation described in Zgodic et al.

\item[\code{sigma\_init}] An initial value for the residual variance parameter. Default is NULL which corresponds to the sample variance of Y.
\end{ldescription}
\end{Arguments}
%
\begin{Value}
A list of the output of the lmmprobe function, including

\code{beta\_hat, beta\_hat\_var} MAP estimates of the posterior expectation (beta\_hat) and variance (beta\_hat\_var) of the prior mean (\eqn{\beta}{}) of the regression coefficients assuming \eqn{\gamma=1}{},

\code{gamma} the posterior expectation of the latent \eqn{\gamma}{} variables,

\code{preds} predictions of \eqn{Y}{},

\code{PI\_lower, PI\_upper} lower and upper prediction intervals for the predictions,

\code{sigma2\_est} MAP estimate of the residual variance,

\code{random\_var} MAP estimate of the random effect(s) variance,

\code{random\_intercept} estimated random intercept terms,

\code{random\_slope} estimated random slope terms, if applicable.
\end{Value}
%
\begin{References}
Zgodic A, Bay R, Zhang J, Olejua P, and McLain AC (2023). Sparse high-dimensional linear mixed modeling with a partitioned empirical Bayes ECM algorithm. arXiv preprint arXiv:2310.12285.
\end{References}
%
\begin{Examples}
\begin{ExampleCode}
library(lmmprobe)
data(SLE)
Y <- matrix(real_data[, "y"], ncol = 1)
Z <- real_data[, 4:ncol(real_data)]
V <- matrix(real_data[, "id"], ncol = 1)
ID_data <- as.numeric(as.character(real_data$id))
full_res <- lmmprobe(Y = Y, Z = Z, V = V, ID_data = ID_data)
\end{ExampleCode}
\end{Examples}
\HeaderA{SLE}{Systemic Lupus Erythematosus (SLE) Data}{SLE}
\aliasA{real\_data}{SLE}{real.Rul.data}
%
\begin{Description}
A dataset containing data from a study on Systemic Lupus Erythematosus (SLE).
Loading this dataset creates an object named `real\_data`.
\end{Description}
%
\begin{Usage}
\begin{verbatim}
data(SLE)
\end{verbatim}
\end{Usage}
%
\begin{Format}
A data frame with observations on the following variables:
\begin{description}

\item[y] Response vector.
\item[id] Subject ID.
\item[time] Time of observation.
\item[...] Other predictors.

\end{description}

\end{Format}
%
\begin{Source}
Simulated or real data provided with the package.
\end{Source}
\printindex{}
\end{document}
